\chapter{巴拿赫空间与基本定理}
\label{ch:b3:banach}

泛函分析研究无穷维赋范空间，
通过定义在其上的算子与泛函进行研究。其
奠基性发现是：借助贝尔定理，\emph{完备性}迫使
强一致性成立：逐点有界的算子族
在范数意义下有界（巴拿赫–斯坦豪斯定理），连续双射
的逆仍连续（开映射定理），图像
的闭性可检验连续性（闭图像定理）。另一根支柱
\emph{哈恩–巴拿赫}定理不需要任何完备性——仅用
佐恩引理——并保证\omterm{def:b3:banach:operator}{对偶空间}足够丰富，
足以分辨每一个向量。本章证明上述四个定理，并
在经典数列空间 $\ell^p$ 上检验它们，
通过具体的对偶计算，以及一个真正出人意料的应用加以检验：
即存在连续的 $2\pi$ 周期函数，
其傅里叶级数在某一点\emph{发散}——以否定方式解决了
本科二年级遗留的一个问题。

全文中 $E, F$ 是定义在 $K = \R$ 或 $\C$ 上的赋范空间；
\emph{巴拿赫空间}\index{巴拿赫空间}指
完备的赋范空间。

\section{有界算子；数列空间}

\begin{definition}\label{def:b3:banach:operator}
$\mathcal L(E, F)$ 表示全体\emph{有界}
（即连续，见本科二年级）线性映射，并赋以\emph{算子
范数}\index{算子范数} $\vertiii T = \sup_{\norm x \leq
1}\norm{Tx}$；它满足次可乘性：$\vertiii{ST} \leq
\vertiii S\,\vertiii T$。\emph{对偶空间}\index{对偶空间}为
$E' = \mathcal L(E, K)$。
\end{definition}

\begin{proposition}\label{prop:b3:banach:llcomplete}
若 $F$ 为巴拿赫空间，则 $\mathcal L(E, F)$ 亦是；特别地，
$E'$ 总是巴拿赫空间。
\end{proposition}

\begin{proof}
设 $(T_n)$ 关于 $\vertiii\cdot$ 为柯西列。对每个 $x$，
$\norm{T_nx - T_mx} \leq \vertiii{T_n - T_m}\norm x$，故 $(T_nx)$
为 $F$ 中的柯西列，因而收敛，记其极限为 $Tx$。$T$ 是
线性的（由线性等式取极限得出）；对
$\norm{T_nx - T_mx} \leq \varepsilon\norm x$（$n, m \geq N$）取极限，
得 $\norm{T_nx - Tx} \leq \varepsilon\norm x$，故 $T_n \to T$
依\omterm{def:b3:banach:operator}{算子范数}，且 $\vertiii T \leq \vertiii{T_N} +
\varepsilon < \infty$。
\end{proof}

\begin{definition}\label{def:b3:banach:lp}
经典数列空间\index{数列空间 $\ell^p$}
（定义在 $K$ 上，以 $\N$ 为指标）：
\[
\ell^p = \Bigl\{x = (x_n) : \norm x_p = \Bigl(\sum_n
\abs{x_n}^p\Bigr)^{1/p} < \infty\Bigr\}\quad (1 \leq p <
\infty),
\]
\[
\ell^\infty = \{x : \norm x_\infty = \sup\abs{x_n} < \infty\},
\]
而 $c_0 = \{x : x_n \to 0\}$ 赋以 $\norm\cdot_\infty$。
$\norm\cdot_p$ 是范数，由闵可夫斯基不等式推出，
离散情形的证明与 \cref{ch:b3:lp} 中完全相同（或
由本科二年级的有限维不等式求和）。全部
都是巴拿赫空间，且 $c_0$ 是
$\ell^\infty$ 的闭子空间（见 \cref{exo:b3:banach:3}）。
\end{definition}

\begin{omfigure}
\begin{tikzpicture}[scale=1.5]
  \draw[->] (-1.45,0) -- (1.45,0) node[right,font=\small]
    {$x_1$};
  \draw[->] (0,-1.45) -- (0,1.45) node[above,font=\small]
    {$x_2$};
  % p = 1 diamond
  \draw[omDef, thick] (1,0) -- (0,1) -- (-1,0) -- (0,-1)
    -- cycle;
  % p = 2 circle
  \draw[omProp, thick] (0,0) circle (1);
  % p = 4 superellipse |x|^4 + |y|^4 = 1
  \draw[omThm, thick, smooth]
    plot[domain=0:90, samples=46]
    ({cos(\x)/(cos(\x)^4 + sin(\x)^4)^0.25},
     {sin(\x)/(cos(\x)^4 + sin(\x)^4)^0.25});
  \draw[omThm, thick, smooth]
    plot[domain=90:180, samples=46]
    ({cos(\x)/(cos(\x)^4 + sin(\x)^4)^0.25},
     {sin(\x)/(cos(\x)^4 + sin(\x)^4)^0.25});
  \draw[omThm, thick, smooth]
    plot[domain=180:270, samples=46]
    ({cos(\x)/(cos(\x)^4 + sin(\x)^4)^0.25},
     {sin(\x)/(cos(\x)^4 + sin(\x)^4)^0.25});
  \draw[omThm, thick, smooth]
    plot[domain=270:360, samples=46]
    ({cos(\x)/(cos(\x)^4 + sin(\x)^4)^0.25},
     {sin(\x)/(cos(\x)^4 + sin(\x)^4)^0.25});
  % p = infty square
  \draw[black!60, thick] (-1,-1) rectangle (1,1);
  \node[omDef, font=\small] at (0.40,0.40) {$p{=}1$};
  \node[omProp, font=\small] at (-0.52,0.75) {$p{=}2$};
  \node[omThm, font=\small] at (0.90,-0.52) {$p{=}4$};
  \node[black!60, font=\small] at (-1.62,1.15)
    {$p{=}\infty$};
  \draw[black!60, ->] (-1.32,1.11) -- (-1.02,1.0);
\end{tikzpicture}

{\small 平面上各 $p$ 范数的单位球，彼此嵌套：
随 $p$ 由 $1$（菱形）增大，经 $2$（圆盘）、$4$（超椭圆）
至 $\infty$（正方形）。每个球的凸性
\emph{正是}闵可夫斯基不等式；在 $p = 1$ 与
$p = \infty$ 两处的角点，使严格凸性、
最佳逼近的唯一性，以及
\cref{exo:b3:lp:12} 中的等号情形同时退化。}
\end{omfigure}

\begin{proposition}[诺伊曼级数]\label{prop:b3:banach:neumann}
设 $E$ 为巴拿赫空间，且 $T \in \mathcal L(E) = \mathcal L(E,E)$
满足 $\vertiii T < 1$。则 $I - T$ 在 $\mathcal
L(E)$ 中可逆，且 $(I - T)^{-1} = \sum_{n\geq0}T^n$（按
\omterm{def:b3:banach:operator}{算子范数}收敛）。因此可逆算子的全体构成
\omterm{def:b3:topology:topology}{开集}，且求逆运算在其上连续。\index{诺伊曼级数}
\end{proposition}

\begin{proof}
级数绝对收敛（$\vertiii{T^n} \leq \vertiii
T^n$，几何级数控制），在巴拿赫空间 $\mathcal L(E)$ 中
（\cref{prop:b3:banach:llcomplete}；
\cref{exo:b3:complete:1}(b)）。裂项相消，$(I - T)\sum_{n \leq
N}T^n = I - T^{N+1} \to I$，另一侧同理。
关于开性：若 $S$ 可逆，且 $\vertiii{H} <
1/\vertiii{S^{-1}}$，则 $S + H = S(I + S^{-1}H)$，其中
$\vertiii{S^{-1}H} < 1$，故可逆。求逆运算的连续性：
由级数表达式给出 $\vertiii{(S+H)^{-1} - S^{-1}} =
O(\vertiii H)$，局部成立。
\end{proof}

\begin{example}[沃尔泰拉方程，用诺伊曼级数求解]
\label{ex:b3:banach:volterraexample}
在 $E = \mathcal C(\intcc01)$ 上，考虑积分方程
\[
u(x) = 1 + \lambda\int_0^xu(t)\,\dd t,
\qquad\text{即}\qquad u = \mathbf 1 + \lambda Tu,
\quad (Tu)(x) = \int_0^xu .
\]
此处 $\vertiii T \leq 1$，故当 $\abs\lambda < 1$ 时，\omterm{prop:b3:banach:neumann}{诺伊曼
级数}可直接应用：$u = (I - \lambda T)^{-1}\mathbf 1 =
\sum_n\lambda^nT^n\mathbf 1$。计算可得 $T^n\mathbf 1 =
\frac{x^n}{n!}$，于是
\[
u(x) = \sum_{n\geq0}\frac{(\lambda x)^n}{n!} =
\eu^{\lambda x} ,
\]
求导即可验证。更好的是 $\vertiii{T^n} \leq
\frac1{n!}$（迭代核按阶乘速度缩小），因此
$\sum\lambda^nT^n$ 对\emph{每个} $\lambda$ 都收敛——
算子 $I - \lambda T$ 对所有 $\lambda
\in \C$ 均可逆，尽管 $\vertiii{\lambda T} \geq 1$ 终将成立：
关键在于幂的谱衰减，而非
一阶范数。沃尔泰拉算子内蕴这种阶乘型衰减
（\cref{exo:b3:complete:4}(a) 正是利用了这一点），因此初值问题
从不会受边值问题那类
共振现象的困扰
（见\cref{ch:b3:spectral}）。
\end{example}

\section{哈恩--巴拿赫}

\begin{theorem}[哈恩--巴拿赫定理，解析形式]
\label{thm:b3:banach:hahnbanach}
设 $E$ 为\emph{实}向量空间，$p \colon E \to \R$
为\emph{次线性}（$p(x + y) \leq p(x) + p(y)$ 且 $p(tx) =
tp(x)$，$t \geq 0$），$F \subseteq E$ 为子空间，$f
\colon F \to \R$ 为线性映射，且 $f \leq p$ 于 $F$ 上。则 $f$ 可延拓
为线性的 $\tilde f \colon E \to \R$，且 $\tilde f \leq p$
于 $E$ 上。\index{哈恩--巴拿赫定理}
\end{theorem}

\begin{proof}
\emph{一步延拓。}设 $x_0 \notin F$，将 $f$ 延拓到
$F \oplus \R x_0$，适当选取 $\alpha = \tilde f(x_0)$
即需对所有 $y \in F$，$t > 0$，
\[
f(y) + t\alpha \leq p(y + tx_0)
\quad\text{且}\quad
f(y) - t\alpha \leq p(y - tx_0),
\]
两式除以 $t$（利用次线性）后化为
\[
\sup_{v \in F}\ \bigl[f(v) - p(v - x_0)\bigr]
\;\leq\; \alpha \;\leq\;
\inf_{u \in F}\ \bigl[p(u + x_0) - f(u)\bigr].
\]
满足要求的 $\alpha$ 存在，当且仅当每一左端项都 $\leq$
每一右端项。事实上 $f(v) + f(u) = f(u + v) \leq p(u + v) \leq
p(u + x_0) + p(v - x_0)$，即 $f(v) - p(v - x_0) \leq p(u +
x_0) - f(u)$。

\emph{佐恩引理。}将 $f$ 的受 $p$ 控制的延拓（序对：子空间，泛函）
按延拓关系排序；一条链的并是其上界；
极大元须定义在全部
$E$ 上，否则一步延拓将与极大性矛盾。
\end{proof}

\begin{corollary}\label{cor:b3:banach:hbnormed}
设 $E$ 为赋范空间（$K = \R$ 或 $\C$）。
\begin{enumerate}
  \item 每个 $f \in F'$（$F$ 为子空间）都可延拓为 $\tilde f
        \in E'$，使得 $\norm{\tilde f}_{E'} = \norm f_{F'}$。
  \item 对每个 $x \neq 0$，存在 $f \in E'$，使 $\norm f
        = 1$，且 $f(x) = \norm x$。特别地，$E'$ 分离
        $E$ 中的点，且 $\norm x = \sup_{\norm f \leq
        1}\abs{f(x)}$。
  \item 对闭子空间 $F$ 及 $x \notin F$，存在 $f
        \in E'$，在 $F$ 上为零，使 $f(x) = d(x, F)$，
        且 $\norm f \leq 1$。
\end{enumerate}
\end{corollary}

\begin{proof}
(1) 实情形：应用\cref{thm:b3:banach:hahnbanach}，取 $p(x)
= \norm f_{F'}\,\norm x$（次线性）；延拓满足
$\pm\tilde f(x) = \tilde f(\pm x) \leq p(x)$，因此
$\norm{\tilde f} \leq \norm f$，而 $\geq$ 则由限制直接得到。
复情形：令 $u = \operatorname{Re}f$，它是一个实泛函
满足 $\abs u \leq \norm f\norm\cdot$；注意 $f(x) = u(x) -
\iu\,u(\iu x)$（在实部与虚部上验证：
$\operatorname{Im}f(x) = -\operatorname{Re}f(\iu x)$）。将
$u$ 以相同的界作实线性延拓，并令 $\tilde f(x) =
\tilde u(x) - \iu\tilde u(\iu x)$：$\C$ 线性（直接检验
乘以 $\iu$ 的情形），它延拓了 $f$；关于范数，对给定的 $x$
写 $\tilde f(x) = r\eu^{\iu\theta}$，则 $\abs{\tilde f(x)}
= \tilde f(\eu^{-\iu\theta}x) = \tilde u(\eu^{-\iu\theta}x)
\leq \norm f\,\norm x$。

(2) 在 $F = Kx$ 上定义 $f(tx) = t\norm x$，并使其范数为 $1$（于 $F$ 上）；
由 (1) 延拓。对偶公式中，$\leq$ 显然，$\geq$ 由
此 $f$ 给出。

(3) 在 $F \oplus Kx$ 上定义 $f(y + tx) = t\,d(x, F)$；
当 $t \neq 0$ 时，$\norm{y + tx} = \abs t\,\norm{x + y/t} \geq \abs
t\,d(x, F) = \abs{f(y + tx)}$，于是在该
子空间上 $\norm f \leq 1$；由 (1) 延拓。
\end{proof}

\begin{remark}\label{rem:b3:banach:bidual}
由 (2)，典范映射 $J \colon E \to E''$，$J(x)(f) =
f(x)$ 是等距映射（见\cref{exo:b3:banach:10}）。每个赋范空间都
嵌入其二次对偶中。若 $J$ 为满射，则称该空间
为\emph{自反}\index{自反}；周末问题
表明 $\ell^p$（$1 < p < \infty$）\omterm{rem:b3:banach:bidual}{自反}，而
$\ell^1$ 则不然。
\end{remark}

\section{贝尔三部曲}

\begin{theorem}[巴拿赫--斯坦豪斯，一致有界性]
\label{thm:b3:banach:banachsteinhaus}
设 $E$ 为\emph{巴拿赫空间}，$F$ 为赋范空间，$(T_i)_{i\in
I} \subseteq \mathcal L(E, F)$ 的一族算子，满足 $\sup_i
\norm{T_ix} < \infty$ 对每个 $x \in E$ 成立。则 $\sup_i
\vertiii{T_i} < \infty$。\index{巴拿赫--斯坦豪斯定理}
\end{theorem}

\begin{proof}
集合 $F_n = \{x : \sup_i\norm{T_ix} \leq n\}$ 皆为闭集
（闭球原像的交），且覆盖 $E$。
由\cref{thm:b3:complete:baire}，存在 $n_0$ 与一个球
$B(x_0, r) \subseteq F_{n_0}$。对 $\norm z < r$，$\norm{T_iz}
\leq \norm{T_i(x_0 + z)} + \norm{T_ix_0} \leq 2n_0$，于是
$\vertiii{T_i} \leq 2n_0/r$，对每个 $i$ 成立。
\end{proof}

\begin{corollary}\label{cor:b3:banach:bslimits}
若 $E$ 为巴拿赫空间，且 $T_n \in \mathcal L(E,F)$ 以\emph{逐点}
方式收敛（$T_nx \to Tx$ 对每个 $x$ 成立），则 $\sup_n
\vertiii{T_n} < \infty$，$T \in \mathcal L(E, F)$，且
$\vertiii T \leq \liminf \vertiii{T_n}$。
\end{corollary}

\begin{proof}
收敛数列有界，此即逐点有界；
由巴拿赫--斯坦豪斯定理，\omterm{def:b3:banach:operator}{算子范数}被某个 $M$ 所界；于是 $\norm{Tx}
= \lim\norm{T_nx} \leq M\norm x$（$T$ 作为逐点极限是线性的），
并对 $\liminf$ 作用于
$\norm{T_nx} \leq \vertiii{T_n}\norm x$ 两端，即得更精确的界。
\end{proof}

\begin{theorem}[发散的傅里叶级数]
\label{thm:b3:banach:fourierdiverge}
存在连续的 $2\pi$ 周期函数 $f$，其
傅里叶级数在 $0$ 处发散：$\sup_N\abs{S_N(f)(0)} =
\infty$。事实上，这样的 $f$ 构成 $\mathcal
C(S^1)$ 的一个稠密子集。\index{傅里叶级数，发散}
\end{theorem}

\begin{proof}
在 $E = (\mathcal C(S^1), \norm\cdot_\infty)$ 中考虑，它是巴拿赫
空间，并考虑泛函 $\Lambda_N(f) = S_N(f)(0) =
\frac1{2\pi}\int_{-\pi}^{\pi}f(t)\,D_N(t)\,\dd t$（狄利克雷
核，本科二年级）。每个 $\Lambda_N$ 连续，且
\[
\norm{\Lambda_N} = \frac1{2\pi}\int_{-\pi}^\pi\abs{D_N(t)}\,\dd
t \;=\; L_N .
\]
（$\leq$ 显然；$\geq$：取连续的 $f$，$\norm f_\infty
\leq 1$，逼近 $\operatorname{sign}D_N$——其符号
只有有限个跳跃点；在每个跳跃点附近长为
$\varepsilon$ 的区间上作光滑化，积分改变量为 $O(N\varepsilon)$）。
\emph{勒贝格常数} $L_N$ 趋于无穷：
\[
L_N = \frac1{2\pi}\int_{-\pi}^{\pi}
\frac{\abs{\sin\bigl((N{+}\tfrac12)t\bigr)}}{\abs{\sin(t/2)}}
\,\dd t
\geq \frac{2}{\pi}\int_0^{\pi}
\frac{\abs{\sin\bigl((N{+}\tfrac12)t\bigr)}}{t}\,\dd t
= \frac{2}{\pi}\int_0^{(N+\frac12)\pi}\frac{\abs{\sin
u}}{u}\,\dd u
\]
（利用 $\abs{\sin(t/2)} \leq t/2$，其在 $[0, \pi]$ 上成立，再
作代换 $u = (N + \tfrac12)t$）。将积分分拆为各个拱：
\[
\int_{(k-1)\pi}^{k\pi}\frac{\abs{\sin u}}u\,\dd u
\geq \frac1{k\pi}\int_{(k-1)\pi}^{k\pi}\abs{\sin u}\,\dd u =
\frac{2}{k\pi},
\qquad\text{于是}\qquad
L_N \geq \frac{4}{\pi^2}\sum_{k=1}^N\frac1k
\xrightarrow[N\to\infty]{} \infty .
\]
若每个连续的 $f$ 都有 $\sup_N
\abs{\Lambda_N(f)} < \infty$，则由巴拿赫--斯坦豪斯定理可推出
$\sup_N\norm{\Lambda_N} < \infty$，矛盾。于是存在某个 $f$
——事实上，这样的 $f$ 构成一个非贫乏的\omterm{def:b3:topology:interior}{稠密集}（它是
$\bigcup_M\{f: \sup_N\abs{\Lambda_Nf}\leq M\}$ 的补集；该并集为可数个闭集之并，
由上述论证，它在任何球中都无内点，
故为\omterm{rem:b3:complete:meagre}{贫乏集}）——这样的函数满足 $\sup_N\abs{S_Nf(0)} =
\infty$。
\end{proof}

\begin{theorem}[开映射定理]\label{thm:b3:banach:openmapping}
设 $E, F$ 为巴拿赫空间，$T \in \mathcal L(E, F)$
为\emph{满射}。则 $T$ 为开映射：$T(B_E(0,1)) \supseteq
B_F(0, c)$，对某个 $c > 0$ 成立。因此，巴拿赫空间之间的有界双射
算子存在有界
逆。\index{开映射定理}
\end{theorem}

\begin{proof}
记 $B = B_E(0,1)$。由满射性得 $F = \bigcup_n
\overline{T(nB)} = \bigcup_n n\,\overline{T(B)}$；由贝尔
定理（\cref{thm:b3:complete:baire}）可知
$\overline{T(B)}$ 有内点，即存在 $B_F(y_0, 4c) \subseteq
\overline{T(B)}$。将中心移至 $0$：对 $\norm y < 4c$，
$y_0 + y$ 与 $y_0$ 均为像 $Tu_k$、$Tv_k$ 的极限，其中
$u_k, v_k \in B$，于是 $y = \lim T(u_k - v_k)$，且 $u_k - v_k
\in 2B$，因而 $B_F(0, 4c) \subseteq \overline{T(2B)}$，即
$B_F(0, 2c) \subseteq \overline{T(B)}$。

\emph{去掉\omterm{def:b3:topology:interior}{闭包}}（此处用到 $E$ 的完备性）：
设 $\norm y < c$。取 $x_1 \in \frac12 B$，使 $\norm{y -
Tx_1} < c/2$（由 $B_F(0,2c) \subseteq \overline{T(B)}$ 缩放 $\frac12$ 得）；
归纳地取 $x_k \in 2^{-k}B$，使 $\norm{y -
T(x_1 + \dots + x_k)} < c\,2^{-k}$。级数 $\sum x_k$
在巴拿赫空间 $E$ 中绝对收敛，其和为 $x \in B$（范数 $<
\sum 2^{-k} = 1$），由连续性知 $Tx = y$，从而 $B_F(0, c)
\subseteq T(B)$。任意\omterm{def:b3:topology:topology}{开集}上 $T$ 的开性由平移与伸缩得到；
对推论而言，$T$ 的开性意味着
$T^{-1}$ 连续。
\end{proof}

\begin{corollary}[等价范数]\label{cor:b3:banach:equivnorms}
若向量空间关于两个可比较的范数
（$\norm\cdot_a \leq C\norm\cdot_b$）均完备，则这两个范数等价。
\end{corollary}

\begin{proof}
恒等映射 $(E, \norm\cdot_b) \to (E, \norm\cdot_a)$
是巴拿赫空间之间的有界双射，其逆
有界。
\end{proof}

\begin{theorem}[闭图像定理]\label{thm:b3:banach:closedgraph}
设 $E, F$ 为巴拿赫空间，$T \colon E \to F$ 为线性映射。若
图像 $\Gamma = \{(x, Tx)\}$ 在 $E \times F$ 中是闭的（即
$x_n \to x$ 且 $Tx_n \to y$ 蕴含 $y = Tx$），则 $T$
有界。\index{闭图像定理}
\end{theorem}

\begin{proof}
$E \times F$ 在范数 $\norm{(x,y)} = \norm x + \norm y$ 下为巴拿赫空间；
$\Gamma$ 作为闭子空间亦为巴拿赫空间。投影
$\pi_E\colon \Gamma \to E$ 是有界双射，因此其
逆 $x \mapsto (x, Tx)$ 有界
（由\cref{thm:b3:banach:openmapping}）：$\norm{Tx} \leq \norm{(x,
Tx)} \leq C\norm x$。
\end{proof}

\begin{method}\label{met:b3:banach:trilogy}
何时使用哪个定理。\emph{哈恩--巴拿赫}：用于
构造具有指定性质的泛函（使某一向量取到范数，
在子空间上取零、由子空间延拓）——
无需完备性。\emph{巴拿赫--施泰因豪斯}：用于将
逐点信息转化为一致界——通常用于证明
某极限运算是连续的，或（用逆否命题）证明
某个元素\emph{发散}，如傅里叶级数的情形。
\emph{开映射 / 闭图像}：无需额外假设，可由代数双射或
由图像的闭性质得到连续性——
典型用途为比较两个完备范数，或证明
自动连续性。三条贝尔型定理均要求
\emph{源空间}完备；否则有反例
（见\cref{exo:b3:banach:7}）。
\end{method}

\section{具体的对偶空间}

\begin{theorem}\label{thm:b3:banach:duals}
等距地成立：$(c_0)' \cong \ell^1$ 与 $(\ell^1)' \cong
\ell^\infty$，经由配对 $\langle x, y\rangle = \sum_n
x_ny_n$。\index{数列空间的对偶}
\end{theorem}

\begin{proof}
证明 $(c_0)' \cong \ell^1$；第二个同构
见\cref{exo:b3:banach:5}。对 $y \in \ell^1$，定义 $\Lambda_y(x)
= \sum x_ny_n$（$x \in c_0$）：绝对收敛，且
$\abs{\Lambda_y(x)} \leq \norm x_\infty\norm y_1$，故
$\norm{\Lambda_y} \leq \norm y_1$。反之，设 $\Lambda \in
(c_0)'$；令 $y_n = \Lambda(e_n)$（$e_n$ 为单位数列）。
对任意 $N$，取检验元 $x^{(N)} = \sum_{n \leq N}
\operatorname{sign}(\overline{y_n})\,e_n \in c_0$（范数 $\leq
1$；若为复数，则改用模为 1 的因子 $\bar y_n/\abs
{y_n}$）：$\Lambda(x^{(N)}) = \sum_{n\leq N}\abs{y_n} \leq
\norm\Lambda$。故 $y \in \ell^1$，且 $\norm y_1 \leq
\norm\Lambda$。最后，$\Lambda = \Lambda_y$：二者在
$e_n$ 上一致，从而在有限数列上一致；有限数列在 $c_0$ 中稠密（截断）：
$\norm{x - \sum_{n \leq N}x_ne_n}_\infty = \sup_{n > N}\abs{x_n}
\to 0$，正是因为 $x_n \to 0$）；连续泛函
在\omterm{def:b3:topology:interior}{稠密集}上一致则相等。该对应是线性的，
双射且等距（$\norm{\Lambda_y} = \norm y_1$ 由
两个不等式得出）。
\end{proof}

\section{练习}

\begin{exercise}[$\star$]\label{exo:b3:banach:1}
计算下列\omterm{def:b3:banach:operator}{算子范数}：
(a) 移位算子 $S(x_1, x_2, \dots) = (0, x_1, x_2, \dots)$ 与
$S^*(x_1, x_2, \dots) = (x_2, x_3, \dots)$（均作用于 $\ell^2$）；
(b) 乘法算子 $M_a x = (a_nx_n)$（作用于 $\ell^2$），
其中 $a \in \ell^\infty$；
(c) 泛函 $\Lambda(f) = \int_0^{1/2}f -
\int_{1/2}^1f$，定义于 $\mathcal C(\intcc01)$ 上——证明
$\norm\Lambda = 1$，并说明该范数\emph{不被取到}。
\end{exercise}

\begin{exercise}[$\star$]\label{exo:b3:banach:2}
设 $E$ 为巴拿赫空间，$T \in \mathcal L(E)$ 可逆，且 $S$
满足 $\vertiii{S - T} < 1/\vertiii{T^{-1}}$，证明 $S$ 可逆，
并估计 $\vertiii{S^{-1} - T^{-1}}$。应用：
若线性方程组 $Tx = b$ 在 $T$ 可逆时可解，则
$T$ 的充分小扰动使其仍唯一可解，
并给出解的改变量的定量估计。
\end{exercise}

\begin{exercise}[$\star\star$]\label{exo:b3:banach:3}
(a) 证明 $\ell^1$、$\ell^\infty$ 与 $c_0$ 都是巴拿赫空间，
并证明 $c_0$ 是 $\ell^\infty$ 中
有限序列所成空间的\omterm{def:b3:topology:interior}{闭包}。
(b) 证明 $\ell^p \subseteq \ell^q$，且 $\norm\cdot_q \leq
\norm\cdot_p$（$1 \leq p \leq q \leq \infty$），并证明
这一包含是严格的。
\end{exercise}

\begin{exercise}[$\star\star$]\label{exo:b3:banach:4}
设 $F \subseteq E$ 为闭子空间，$x \notin F$。
利用\cref{cor:b3:banach:hbnormed}，证明对偶公式
\[
d(x, F) = \max\bigl\{\abs{f(x)} : f \in E',\ \norm f \leq 1,\
f\restriction_F = 0\bigr\}
\]
（注意这里取的是\emph{最大值}）。由此推出 $F = \bigcap\{\ker f : f
\in E',\ f\restriction_F = 0\}$：闭子空间恰好是
若干泛函之核的交集。
\end{exercise}

\begin{exercise}[$\star\star$]\label{exo:b3:banach:5}
证明 $(\ell^1)' \cong \ell^\infty$ 为等距同构，仿照
\cref{thm:b3:banach:duals}的思路（有限序列
在 $\ell^1$ 中稠密）。
该论证在 $(\ell^\infty)'$ 的情形下何处失效？
\end{exercise}

\begin{exercise}[$\star\star$]\label{exo:b3:banach:6}
设 $E, F, G$ 为赋范空间，$E$ 为巴拿赫空间，且 $B \colon E \times
F \to G$ 为双线性映射，关于每个变量分别连续。证明
$B$ 是（联合）连续的：$\norm{B(x,y)} \leq
C\norm x\norm y$。\emph{（对族
$(B(\cdot, y))_{\norm y \leq 1}$ 应用巴拿赫–斯坦因豪斯定理。）}
\end{exercise}

\begin{exercise}[$\star\star$]\label{exo:b3:banach:7}
(a) 在 $E = \mathcal C(\intcc01)$ 上比较 $\norm\cdot_\infty$
与 $\norm\cdot_1$：恒等映射 $(E, \norm\cdot_\infty) \to
(E, \norm\cdot_1)$ 是有界双射，但其逆映射
无界。在
\cref{cor:b3:banach:equivnorms}中哪一条假设不成立？
(b) 给出一个从 $\ell^2$ 的稠密子空间
到 $K$ 的不连续线性映射（例如定义在有限序列上），并说明
这为何并不与闭图像定理相矛盾。
\end{exercise}

\begin{exercise}[$\star\star\star$]\label{exo:b3:banach:8}
（赫林格–托普利茨）设 $T \colon \ell^2 \to \ell^2$ 为线性映射
（处处有定义）且\emph{对称}：$\langle Tx,
y\rangle = \langle x, Ty\rangle$ 对所有 $x, y$ 成立，其中 $\langle
x, y \rangle = \sum x_n\bar y_n$。证明 $T$ 有界。
\emph{（闭图像：若 $x_k \to x$ 且 $Tx_k \to z$，则对
任意 $y$ 作检验。）}要点：无界对称算子
——量子力学中的哈密顿量——永远不能
定义在整个空间上。
\end{exercise}

\begin{exercise}[$\star\star\star$]\label{exo:b3:banach:9}
（波利亚的求积定理）对每个 $n$，令 $\Lambda_n(f)
= \sum_{i=0}^{n} w_{i,n}f(x_{i,n})$ 为
$\mathcal C(\intcc01)$ 上的求积公式（$x_{i,n} \in \intcc01$，$w_{i,n} \in
\R$）。证明 $\Lambda_n(f) \to \int_0^1f$ 对\emph{每个}
连续的 $f$ 成立，当且仅当：(i) $\Lambda_n(P) \to \int_0^1P$
对每个多项式 $P$ 成立，并且 (ii) $\sup_n\sum_i\abs{w_{i,n}} <
\infty$。\emph{（计算 $\norm{\Lambda_n}$；用
巴拿赫–斯坦因豪斯定理与魏尔斯特拉斯定理。）}验证具有
\emph{正}权重且对常数精确的求积公式满足
(ii)，这是自动成立的。
\end{exercise}

\begin{exercise}[$\star\star$]\label{exo:b3:banach:10}
证明 $J \colon E \to E''$，$J(x)(f) = f(x)$ 为线性
等距映射（利用\cref{cor:b3:banach:hbnormed}(2)），并证明
当 $\dim E < \infty$ 时它是满射。又证明若 $E'$ 可分，
则 $E$ 亦可分。\emph{（取 $x_n$ 近似达到 $E'$ 稠密序列中元素的范数，
并证明它们的闭线性张成为 $E$，
借助\cref{cor:b3:banach:hbnormed}(3)。）}
\end{exercise}

\begin{exercise}[$\star\star$]\label{exo:b3:banach:11}
（商空间）设 $E$ 为巴拿赫空间，$F \subseteq
E$ 为\emph{闭}子空间。在 $E/F$ 上定义
\[
\norm{\bar x} \;=\; d(x, F) = \inf_{y\in F}\norm{x - y} .
\]
(a) 证明这是 $E/F$ 上良定义的\emph{范数}（
$F$ 的闭性在何处起作用？），并证明投影 $\pi
\colon E \to E/F$ 满足 $\vertiii\pi \leq 1$，并将开
单位球映满开单位球。
(b) 证明 $E/F$ 是完备的。\emph{（利用
\cref{exo:b3:complete:1}(b) 的级数判别法：给定等价类
$\bar x_k$ 满足 $\sum\norm{\bar x_k} < \infty$，将每个提升为
$x_k \in E$，使得 $\norm{x_k} \leq \norm{\bar x_k} + 2^{-k}$，
并在 $E$ 中求和。）}
(c) 计算：对 $E = c$（收敛数列）与 $F =
c_0$，证明 $c/c_0 \cong K$ 为等距同构，映射为 $\bar x \mapsto
\lim_nx_n$。
\end{exercise}

\begin{exercise}[$\star\star$]\label{exo:b3:banach:12}
（有界投影与可补子空间）设 $E$ 为
巴拿赫空间，$P \colon E \to E$ 为线性映射，满足 $P^2 = P$
（代数投影），$V = \operatorname{im}P$，$W =
\ker P$.
(a) 设 $P$ 有界。证明 $V$ 与 $W$ 都是闭的，
且 $E = V \oplus W$，分解式为 $x = Px + (x -
Px)$.
(b) 反之，设 $E = V \oplus W$，且\emph{两者} $V,
W$ 均为闭的，令 $P$ 为到 $V$ 上沿 $W$ 的投影。
证明 $P$ 有界。\emph{（闭图像：若 $x_n \to x$
且 $Px_n \to z$，则 $z \in V$，$x_n - Px_n \to x - z \in
W$，由分解的唯一性可知 $z =
Px$.)}
(c) 推出\emph{等价性}：子空间 $V$ 容许
有界投影当且仅当它是闭的，且具有闭的
代数补——并（不加证明地）指出，存在不具此性质的闭子空间
（$c_0$ 在
$\ell^\infty$ 中是经典例子）：希尔伯特空间，
其中 $V^\perp$ 总是可用，
（见\cref{ch:b3:hilbert}）是例外而非通则。
\end{exercise}

\section{问题：\texorpdfstring{$\ell^p$}%
{lp} 空间的对偶性}

\begin{problem}[{周末问题——$(\ell^p)' = \ell^q$，
自反性，以及 $\ell^\infty$ 的奇异性}]
\label{pb:b3:banach:1}
设 $1 < p < \infty$，$q$ 为其共轭指数，
$\frac1p + \frac1q = 1$。全文所用的配对为 $\langle x,
y\rangle = \sum_n x_ny_n$。

\medskip
\noindent\textbf{第一部分——数列上的赫尔德与闵可夫斯基
不等式。}
\begin{enumerate}
  \item （杨氏不等式）对 $a, b \geq 0$，证明 $ab \leq
        \frac{a^p}p + \frac{b^q}q$，利用 $\log$ 的凹性，或
        通过研究 $t \mapsto \frac{t^p}p + \frac1q - t$。
  \item （赫尔德）推出：$\abs{\langle x, y\rangle} \leq
        \norm x_p\norm y_q$，其中 $x \in \ell^p$，$y \in \ell^q$；
        指出等号成立的情形。
  \item （闵可夫斯基）推出
        $\norm\cdot_p$ 的三角不等式。\emph{（写出 $\abs{x_n + y_n}^p \leq
        \abs{x_n}\,\abs{x_n{+}y_n}^{p-1} +
        \abs{y_n}\,\abs{x_n{+}y_n}^{p-1}$，并对
        每一项应用赫尔德不等式。）}
  \item 证明 $\ell^p$ 完备，且有限数列
        在其中稠密。
\end{enumerate}

\medskip
\noindent\textbf{第二部分——对偶性 $(\ell^p)' =
\ell^q$。}
\begin{enumerate}[resume]
  \item 对 $y \in \ell^q$，证明 $\Lambda_y(x) = \langle
        x, y\rangle$ 定义了 $\Lambda_y \in (\ell^p)'$，且
        $\norm{\Lambda_y} \leq \norm y_q$，并在
        $x_n = \abs{y_n}^{q-1}\operatorname{sign}(y_n)$
        （经适当截断与规范化）后的该序列上检验，
        证明 $\norm{\Lambda_y} = \norm y_q$。
  \item 反之，给定 $\Lambda \in (\ell^p)'$，令 $y_n =
        \Lambda(e_n)$；证明 $y \in \ell^q$，且 $\norm y_q \leq
        \norm\Lambda$ \emph{（对截断序列做检验，如
        第 5 题，令截断长度趋于无穷）}，并得出
        $\Lambda = \Lambda_y$：映射 $y \mapsto \Lambda_y$ 是
        一个等距同构 $\ell^q \to (\ell^p)'$。
  \item 推出 $\ell^p$ 是\emph{\omterm{rem:b3:banach:bidual}{自反}}的，当 $1 < p <
        \infty$ 时：复合两次对偶后，
        $(\ell^p)''$ 的每个元素都来自 $\ell^p$；仔细验证该复合即为
        典范映射 $J$。
\end{enumerate}

\medskip
\noindent\textbf{第三部分——$\ell^1$ 与 $\ell^\infty$ 是
截然不同的两类对象。}
\begin{enumerate}[resume]
  \item 证明 $\ell^p$（$1 \leq p < \infty$）与 $c_0$ 均
        可分，但 $\ell^\infty$ 不可分。\emph{（
        $\N$ 的子集所对应的不可数多个示性数列
        两两距离为 $1$。）}
  \item 由\cref{exo:b3:banach:10}可知 $(\ell^1)'
        \cong \ell^\infty$，但 $(\ell^\infty)' \not\cong
        \ell^1$：$\ell^1$ \emph{不}是\omterm{rem:b3:banach:bidual}{自反}的。
        \emph{（若 $(\ell^\infty)'$ 为 $\ell^1$，则它将是
        可分的，从而迫使 $\ell^\infty$ 可分。）}
  \item （显式的巴拿赫极限）在 $\ell^\infty_\R$ 上，令
        $p(x) = \limsup_n \frac{x_1 + \dots + x_n}{n}$。证明
        $p$ 是次线性的，并且在子空间 $c$ 上，
        （即收敛数列全体），$\mathrm{LIM}(x) = \lim x$
        满足 $\mathrm{LIM} \leq p$。由哈恩–巴拿赫定理，延拓为
        $\mathrm{LIM} \colon \ell^\infty_\R \to \R$，并证明：
        $\mathrm{LIM}$ 是正的（$x \geq 0 \Rightarrow
        \mathrm{LIM}(x) \geq 0$），平移不变
        （$\mathrm{LIM}(x_2, x_3, \dots) = \mathrm{LIM}(x)$），
        延拓了极限，并满足 $\liminf x \leq
        \mathrm{LIM}(x)\leq \limsup x$。
  \item 证明：把这样的 $\mathrm{LIM}$ 看作
        $(\ell^\infty)'$ 中的元素时，它\emph{不}具有 $\Lambda_y$ 的形式，
        对任意 $y \in \ell^1$ 均成立；由此再次得出 $(\ell^\infty)'
        \neq \ell^1$。\emph{（在单位序列
        $e_n$ 上求值，再在常数列 $1$ 上求值。）}
  \item 计算 $\mathrm{LIM}$ 在 $(0,1,0,1,\dots)$ 上的取值，并证明
        不存在平移不变的\emph{乘性}
        极限延拓，\emph{（考虑 $x =
        (0,1,0,1,\dots)$ 与 $x\cdot Sx$，其中 $S$ 为
        移位算子）}。
\end{enumerate}

\medskip
\noindent\textbf{第四部分——尾声：\omterm{rem:b3:banach:bidual}{自反}性
为何重要。}
\begin{enumerate}[resume]
  \item 利用\cref{cor:b3:banach:bslimits}及第 6 题，
        证明：$\ell^p$（$1 < p <
        \infty$）中每个有界数列都有子列 $(x^{(k)})$，该子列
        \emph{弱收敛}：$\Lambda(x^{(k)})$ 对每个
        $\Lambda \in (\ell^p)'$ 都收敛。\emph{（在
        可数多个坐标上作对角抽取；在
        $\ell^p$ 中确定弱极限，利用范数的一致有界性与
        赫尔德不等式。）}通过例子（$e_n$ 在 $\ell^1$ 中，与
        $\ell^\infty$ 中适当选取的元素配对）说明此结论
        在 $\ell^1$ 中失效：弱紧性是\omterm{rem:b3:banach:bidual}{自反}空间的
        特权。
\end{enumerate}

\medskip
\noindent\textbf{第五部分——弱\omterm{def:b3:topology:topology}{拓扑}的应用与
舒尔的意外结果。} 称 $x^{(k)} \rightharpoonup x$ 为
赋范空间 $E$ 中的\emph{弱收敛}，即当
$\Lambda(x^{(k)}) \to \Lambda(x)$ 对每个 $\Lambda \in E'$ 成立。
\begin{enumerate}[resume]
  \item 完成这一汇总：证明 $(c_0)' \cong \ell^1$
        等距地，沿用第 5 至 6 题的思路（什么
        替代检验序列？）。组装出链
        $c_0 \to \ell^1 \to \ell^\infty \to \dots$ 的相继对偶，
        并标出\omterm{rem:b3:banach:bidual}{自反}性在何处失效。
  \item 证明：巴拿赫空间中每个弱收敛数列
        都有界：把 $x^{(k)}$ 经
        典范嵌入 $J$ 视为 $E'$ 上的泛函，并
        应用巴拿赫–斯坦豪斯
        （\cref{thm:b3:banach:banachsteinhaus}）——在哪个
        巴拿赫空间上，又为何在那里可以用完备性？
  \item 证明：在 $\ell^p$ 中（$1 < p < \infty$），$x^{(k)}
        \rightharpoonup x$ 当且仅当 $\sup_k\norm{x^{(k)}}_p <
        \infty$，且 $x^{(k)}_n \to x_n$ 对每个坐标
        $n$ 成立 \emph{（一个方向用到巴拿赫–斯坦豪斯，经由
        典范嵌入；另一方向则用有限数列逼近 $y
        \in \ell^q$）}。由此推出 $e_k
        \rightharpoonup 0$ 在 $\ell^2$ 中成立，而 $\norm{e_k}_2 =
        1$：弱极限可能丢失质量。
  \item 证明：范数是弱下半连续的：
        $x^{(k)} \rightharpoonup x$ 蕴含 $\norm x \leq
        \liminf_k\,\norm{x^{(k)}}$ \emph{（取一个使
        $x$ 范数达到的泛函，见\cref{cor:b3:banach:hbnormed}）}。
  \item （$\ell^2$ 中的拉东–里斯性质）证明：在 $\ell^2$ 中，弱
        收敛加上范数收敛，蕴含
        依范数收敛 \emph{（展开 $\norm{x^{(k)} -
        x}_2^2$）}。对去掉范数假设的同一命题，给出
        反例。
  \item （舒尔，第一步）设 $x^{(k)} \rightharpoonup 0$ 在
        $\ell^1$ 中，并用反证法假设
        $\norm{x^{(k)}}_1 \geq \delta > 0$ 沿某个
        子列成立。先证明 $x^{(k)}_n \to 0$ 对
        每个 $n$ 成立（需要哪些泛函？），然后递归地构造
        指标 $k_1 < k_2 < \cdots$ 及整数
        $0 = N_0 < N_1 < N_2 < \cdots$，使得
        $x^{(k_j)}$ 的质量集中于块 $B_j =
        \intoc{N_{j-1}}{N_j}$：
        \[
        \sum_{n \in B_j}\bigl|x^{(k_j)}_n\bigr| \geq
        \norm{x^{(k_j)}}_1 - \frac\delta{10} .
        \]
  \item （舒尔，第二步）定义 $y \in \ell^\infty$，令 $y_n =
        \operatorname{sign}\bigl(x^{(k_j)}_n\bigr)$，当 $n \in
        B_j$ 时。证明 $\langle x^{(k_j)}, y\rangle \geq
        \norm{x^{(k_j)}}_1 - \frac{2\delta}{10} \geq
        \frac{8\delta}{10}$，并导出与
        $x^{(k)} \rightharpoonup 0$ 的矛盾。于是得出\emph{舒尔
        定理}：在 $\ell^1$ 中，弱收敛数列
        依范数收敛。
  \item 由此推出 $(e_k)$ 在 $\ell^1$ 中没有
        弱收敛子列（其唯一候选极限为
        $0$，逐坐标——再应用舒尔定理），从而恢复
        第 13 题中弱紧性失效的结论；并化解表面的悖论：
        在 $\ell^1$ 中，弱收敛与范数收敛
        对\emph{数列}一致，然而
        弱与范数\emph{\omterm{def:b3:topology:topology}{拓扑}}不同，有界集仍
        不是弱序列紧的——
        这并无矛盾，只是\omterm{rem:b3:banach:bidual}{自反}性失效而已。
  \item （综合表）对 $E \in \{c_0,\ \ell^1,\ \ell^p\
        (1{<}p{<}\infty),\ \ell^\infty\}$，列表给出：其\omterm{def:b3:banach:operator}{对偶空间}；
        可分性；\omterm{rem:b3:banach:bidual}{自反}性；有界数列是否
        存在弱收敛子列；以及各空间的
        一条标志性质，
        并由本题在一行内说明其理由。
\end{enumerate}

\medskip
\noindent\textbf{第六部分——补充：最近点、
平均收敛与巴拿赫极限的值}
\begin{enumerate}[resume]
  \item （最近点：\omterm{rem:b3:banach:bidual}{自反}性的应用）设 $F$ 为
        $\ell^p$ 的闭子空间（$1 < p < \infty$），且
        $x \in \ell^p$。证明 $d = \operatorname{dist}(x,
        F)$ \emph{可取到}：由一极小化
        序列抽取一弱收敛子列（第
        13 题），为使弱极限落在 $F$ 内，借助
        哈恩--巴拿赫定理构造一个在 $F$ 上为零、
        而在子空间之外某点不为零的泛函，并由第 17 题作结。
        再证此“特权”并非普遍：在 $c_0$ 中，
        对 $\Lambda(x) = \sum_n2^{-n}x_n$，证明
        $\norm\Lambda = 1$ 在单位球上不被达到，
        建立距离公式
        $\operatorname{dist}(x, \ker\Lambda) =
        \abs{\Lambda(x)}$，并推出：对任意 $x \notin
        \ker\Lambda$，在闭
        超平面 $\ker\Lambda$ 中都没有最近点。
  \item （$\ell^2$ 中的巴拿赫--萨克斯性质）设 $x^{(k)}
        \rightharpoonup 0$ 于 $\ell^2_{\R}$ 中，且
        $\norm{x^{(k)}}_2 \leq C$。构造子列
        $(y_j)$，使 $\abs{\langle y_i, y_j\rangle} \leq
        \frac1j$，对一切 $i < j$ 成立，并由此推出
        \[
        \Bigl\lVert\frac{y_1 + \dots + y_m}m\Bigr\rVert_2^2
        \leq \frac{C^2 + 2}m \longrightarrow 0 :
        \]
        取子列后，切萨罗平均依
        \emph{范数}收敛。对 $(e_k)$ 验证之，其平均的范数为
        $\frac1{\sqrt m}$：弱收敛对数列本身
        并无作用（第 16 题），对平均却化为范数
        收敛。
  \item （巴拿赫极限的值）设 $L$ 为任一巴拿赫
        极限（第 10 题），令 $A_mx = \frac1m(x + Sx + \dots
        + S^{m-1}x)$。证明 $L(A_mx) = L(x)$，且 $\liminf A_mx
        \leq L(x) \leq \limsup A_mx$，由此推出：\emph{所有}
        巴拿赫极限在周期数列上取值一致，其值为
        即一个周期内的平均值——$\frac13$，对数列 $(1, 0, 0, 1,
        0, 0, \dots)$ 的值与第 12 题的
        $\frac12$ 相符。再证明一般情形下一致性不成立：对
        分块数列 $x$，其值为 $1$ 的部分是
        $\intoc{3^{j-1}}{3^j}$（$j$ 为偶数），其余处取值 $0$，
        证明切萨罗平均在 $\leq
        \frac13$ 与 $\geq \frac23$ 之间振荡，并构造两个巴拿赫
        极限 $L_\pm$，使 $L_-(x) \leq \frac13 < \frac23
        \leq L_+(x)$ \emph{（自 $c \oplus \R x$ 出发延拓，
        取可容许的极端值 $\pm$，验证
        $\Lambda(y + tx) = \lim y + t\,p(x)$ 被
        第 10 题中的次线性 $p$ 所控制）}。
\end{enumerate}
\end{problem}
